\section{Configuration and fitted parameters}
\label{app:params}

This appendix records the frozen configuration in enough detail to reconstruct it
exactly. It is a reproducibility artifact. As stated in \S\ref{sec:fitting}, no
strategic interpretation should be placed on individual coefficients: the ask
features are mutually correlated and normalised to different effective ranges, and
two of the decision knobs are inert in the shipping configuration.
Table~\ref{tab:features} and Table~\ref{tab:params} are a pair: the first defines
the \numAskFeats{} ask features, the second gives the weight fitted to each of
them, under the same feature names. The appendix then records the per-coordinate
bounds, the implementation detail of ask scoring and of the two-ply lookahead,
the value-function features, coefficients and provenance gap, the two
endgame decisions that belong to the policy rather than to the driver
(\S\ref{app:params-endgame}), and the specification grammar that names any
configuration on the command line.

\subsection{The ask features}

\begin{table}[ht]
\centering
\caption{The \numAskFeats{} ask features $\varphi_k$ entering the fitted linear
term of \eqref{eq:askutil}, as chosen in \S\ref{sec:ask}. All are normalised to
roughly $[0,1]$ so that the fitted weights are comparable in scale; the weight
fitted to feature $k$ appears in Table~\ref{tab:params} under the same name. $H$
denotes the named card's half-suit and $p = \hat\mu_{c,t}$ the Fast posterior
that the target holds the named card; probabilities written $\Pr[\cdot]$ are Fast
per-card marginals, so every feature is computed from the approximate inference
path rather than from the reference engine. See \S\ref{sec:fitting} for why
individual coefficients should not be read as strategy.}
\label{tab:features}
\small
\begin{tabular}{rlp{0.60\linewidth}}
\toprule
$k$ & Name & Definition \\
\midrule
0  & hit probability     & $p = \hat\mu_{c,t}$, the Fast posterior that the target holds the named card \\
1  & squared hit         & $p^2$; lets the fit curve the trade-off between certainty and value \\
2  & certain hit         & $\mathbf{1}[p > 0.9995]$ \\
3  & own progress        & cards of $H$ in hand, over six \\
4  & team control        & expected fraction of $H$ held by our team \\
5  & lock completion     & $\prod_{d \in H \setminus \{c\}} \Pr[d \text{ is ours}]$: the probability this ask alone completes team ownership of the half-suit \\
6  & continuation        & best posterior on any other missing card of $H$ \\
7  & completion bonus    & $1$ at four or more cards of $H$ in hand, $0.35$ at three \\
8  & reply threat        & $(1-p)\cdot$ the best ask the target could make against us if handed the turn \\
9  & information leak    & $1$ if our team has never publicly revealed an interest in $H$ \\
10 & target hand size    & target's public card count over nine \\
11 & empties the target  & $p$ if the target holds exactly one card, else $0$ \\
12 & repeats our half-suit & $1$ if this repeats our own previous half-suit \\
13 & known team cards    & publicly located cards of $H$ on our team, over six \\
14 & location entropy    & $-p\log_2 p - (1-p)\log_2(1-p)$ \\
15 & team owns $H$       & $\prod_{d \in H}\Pr[d \text{ is ours}]$ before the ask; an independence proxy, not the exact query \eqref{eq:q3} \\
16 & exposure on a miss  & $(1-p)\cdot$ our cards visible in half-suits the target has asked in \\
17 & trailing pressure   & $p$ when behind by two or more half-suits \\
18 & runway              & expected length of the run this ask begins, from the best remaining per-card posteriors \\
19 & leak magnitude      & feature 9 scaled by how much of $H$ we hold \\
\bottomrule
\end{tabular}
\end{table}

\subsection{Ask scoring and lookahead implementation}
\label{app:params-ask}

Two details of \S\ref{sec:ask} and \S\ref{sec:twoply} are recorded here rather
than in the main text.

Ties in \eqref{eq:askutil}, and in the rescoring \eqref{eq:twoply} that follows
it, are broken deterministically: first by the index of the named card, then by
the index of the target seat, so no randomisation enters the choice --- this is
part of what the determinism claimed in \S\ref{sec:policy} rests on.

The branch states of the two-ply refinement are constructed explicitly rather
than sampled. On the post-hit branch the named card is resolved to the asker's
hand and the two affected public hand counts change; on the post-miss branch an
exclusion is added, removing the target from that card's surviving support $M_c$.
Capacities \eqref{eq:capacity} are re-derived on each branch and the Fast
marginals of \S\ref{sec:fast} refitted there, which is why the branch beliefs
carry the same approximation error as $\hat\mu$ itself.

\subsection{The frozen parameter vector}

\begin{table}[ht]
\centering
\caption{The frozen \vfast{} configuration: the \numAskFeats{} fitted ask weights
of Table~\ref{tab:features} (left pair of columns) and the \numKnobs{} fitted decision
knobs (right pair), together the \numParams{} coordinates of the vector searched
in \S\ref{sec:fitting}. Values are those selected at generation
\numSelectedGen{} of round five on validation seed $1357911$ (500 deals, two
rotations, panel \{\vthree{}, lockout, detective, \vtwo{}\}) and recorded in
\texttt{research/v04/runs/selected.json}; the same values are compiled into
\texttt{V04Config} and are used unchanged in every experiment reported in
\S\ref{sec:results}. The two integer coordinates are shown before
rounding: the patience pool is applied as $6$ and the lookahead width $K$ as $6$.
This table is a reproducibility record, not a description of strategy.}
\label{tab:params}
\footnotesize
\begin{tabular}{lrlr}
\toprule
Ask feature & Weight & Decision knob & Value \\
\midrule
hit probability & +11.506 & declare threshold & 0.818 \\
squared hit & +3.295 & locked threshold & 0.797 \\
certain hit & +3.198 & ask floor & 0.333 \\
own progress & +1.688 & patience pool & 6.249 \\
team control & +2.133 & opp.\ card floor & 2.868 \\
lock completion & +4.071 & value weight & 6.043 \\
continuation & +1.468 & linear weight & 0.767 \\
completion bonus & +1.428 & min.\ team prob. & 0.792 \\
reply threat & -3.098 & stopping margin & -0.034 \\
information leak & -0.854 & prior $\theta$ & 0.264 \\
target hand size & -2.022 & prior $\phi$ & 0.133 \\
empties the target & +1.166 & two-ply $K$ & 5.624 \\
repeats our half-suit & +1.270 & chain weight & 3.358 \\
known team cards & +0.914 & threat weight & 2.613 \\
location entropy & -2.653 &  &  \\
team owns $H$ & -0.804 &  &  \\
exposure on a miss & +1.904 &  &  \\
trailing pressure & +0.047 &  &  \\
runway & +1.411 &  &  \\
leak magnitude & -0.999 &  &  \\
\bottomrule
\end{tabular}
\end{table}

The same vector can be supplied to any build of the engine as a single
\texttt{allparams=} string. Concatenate the three lines below with no intervening
whitespace; coordinates are separated by \texttt{|} and appear in the order
\emph{ask weights $0$--$19$, then the \numKnobs{} knobs in the order of
Table~\ref{tab:params}}:

{\footnotesize
\begin{verbatim}
11.506|3.2948|3.1978|1.6881|2.1333|4.0705|1.4679|1.4281|-3.0978|-0.8536|-2.0219|1.166|
1.2697|0.9142|-2.6534|-0.8045|1.904|0.0473|1.4108|-0.999|0.8177|0.7969|0.3325|6.249|
2.8676|6.0432|0.7667|0.7925|-0.0342|0.2638|0.1328|5.6239|3.358|2.6127
\end{verbatim}
}

\subsection{Bounds imposed per coordinate}

The optimiser samples from an unconstrained diagonal Gaussian; the bounds below
are applied when a sampled vector is turned into a configuration, both by the
engine's parameter parser (\filepath{engine/src/factory.hpp}) and by the freeze step
(\filepath{engine/freeze_config.py}), so a coordinate outside its bound is clamped
rather than rejected. The two integer coordinates are rounded to the nearest
integer after clamping.

\begin{table}[ht]
\centering
\caption{Bounds applied to each coordinate of the \numParams{}-coordinate
parameter vector when it is instantiated as a configuration. Coordinates
$0$--$19$ are the ask weights of Table~\ref{tab:features} and are not clamped;
coordinates $20$--$33$ are the decision knobs. Bounds are identical in the engine
parser and in the freeze step, so the vector scored during fitting and the vector
compiled into the shipping binary are subject to the same constraints.}
\label{tab:bounds}
\small
\begin{tabular}{rll}
\toprule
Coordinate & Quantity & Bound \\
\midrule
$0$--$19$ & ask weights $w_0,\dots,w_{19}$ & unbounded \\
20 & declaration threshold & $[0.5,\,0.9999]$ \\
21 & locked-half-suit threshold & $[0.5,\,0.99999]$ \\
22 & ask floor & $[0,\,0.9]$ \\
23 & patience pool & $[0,\,45]$, rounded to an integer \\
24 & opponent-card floor & $[0,\,20]$ \\
25 & value weight $\lambda_{\mathrm{val}}$ & $[0,\,\infty)$ \\
26 & linear weight $\lambda_{\mathrm{lin}}$ & $[0,\,\infty)$ \\
27 & minimum team probability & $[0.05,\,0.99]$ \\
28 & stopping margin $\varepsilon$ & unbounded \\
29 & prior $\theta$ & $[0,\,2]$ \\
30 & prior $\phi$ & $[0,\,1]$ \\
31 & lookahead width $K$ & $[0,\,24]$, rounded to an integer \\
32 & chain weight $\lambda_{\mathrm{chain}}$ & $[0,\,\infty)$ \\
33 & threat weight $\lambda_{\mathrm{threat}}$ & $[0,\,\infty)$ \\
\bottomrule
\end{tabular}
\end{table}

\subsection{Value-function coefficients, and the provenance gap}

The \numValueFeats{} features of $V$ are named in Table~\ref{tab:valuecoef} and
defined as follows: a bias and the current score differential; three control
terms built from $e_H$, the expected fraction of half-suit $H$ our team holds,
namely $\sum_H (2e_H - 1)$, a sharpened version of it and the contested mass
$\sum_H e_H(1-e_H)$; the signed count of half-suits already owned each way;
counts of what is left to play for, from the card differential and the size of
$|U|$ to the near-complete half-suits leaning each way; and side to move with its
interactions with control and with $|U|$. Features are collected from all six
seats at every decision point rather than from the mover alone, because under
mover-only collection the side-to-move feature is constant at $+1$ and its
coefficient --- which the miss branch of \eqref{eq:askutil} and the stopping rule
\eqref{eq:stop} both depend on --- would be unidentified.

The \numValueFeats{} coefficients of the value function $V$ (\S\ref{sec:value})
are not part of the \numParams{}-coordinate vector and are not written by the
freeze step, which edits only the policy parameters. Two distinct vectors
therefore exist in the repository, and they differ numerically:
Table~\ref{tab:valuecoef} prints both. The left column is what
\texttt{V04Config::vw} contains and is what produced every result in
\S\ref{sec:results}. The right column is the ridge fit reported as E14
(\numValueRows{} rows, in-sample $R^2 = \numValueRSq$, RMSE \numValueRmse{}),
stored in \filepath{research/v04/results/E14-valuefit.txt}.

The two vectors come from different runs, and no step in the pipeline copies the
second into the first. The consequence is that the
E14 metrics describe the quality of the artifact's fit and not the quality of the
coefficients that were actually deployed, and that there is no reported
goodness-of-fit figure --- in-sample or held-out --- for the deployed vector. We
record this rather than repair it, because repairing it would change the
configuration that produced every number in this paper.

\begin{table}[ht]
\centering
\caption{The \numValueFeats{} linear coefficients of the value function $V$.
``Compiled'' is the vector in \texttt{V04Config::vw} used by \vfast{} and
\vblock{} in every experiment reported here; ``E14 artifact'' is the separate
ridge fit of \texttt{research/v04/results/E14-valuefit.txt}, obtained on 250
self-play deals at seed $31415$. The two vectors differ because the freeze step
writes only the \numParams{} policy parameters and never the value coefficients.}
\label{tab:valuecoef}
\small
\begin{tabular}{rlrr}
\toprule
$j$ & Feature & Compiled & E14 artifact \\
\midrule
0  & bias                          & $+0.001242$ & $+0.007023$ \\
1  & score differential            & $+0.888965$ & $+0.909779$ \\
2  & expected control              & $+0.421266$ & $+0.115403$ \\
3  & sharpened control             & $-0.145791$ & $-0.249275$ \\
4  & locked differential           & $+0.225573$ & $+0.011511$ \\
5  & side to move                  & $+0.022896$ & $+0.026911$ \\
6  & card differential             & $+0.422207$ & $+0.114251$ \\
7  & unresolved pool               & $+0.007678$ & $-0.079727$ \\
8  & live half-suits               & $+0.005904$ & $+0.035886$ \\
9  & side to move $\times$ control & $-0.000997$ & $-0.003278$ \\
10 & own hand size                 & $-0.006601$ & $-0.011850$ \\
11 & smallest friendly hand        & $-0.007472$ & $+0.062496$ \\
12 & our near-complete half-suits  & $-0.022484$ & $+0.005226$ \\
13 & their near-complete half-suits& $-0.025189$ & $-0.034832$ \\
14 & contested mass                & $+0.080416$ & $+0.027557$ \\
15 & side to move $\times$ unresolved & $-0.021409$ & $-0.019608$ \\
\bottomrule
\end{tabular}
\end{table}

Either vector can be supplied at the command line as a \texttt{vweights=} string,
in the coordinate order of Table~\ref{tab:valuecoef}. The compiled vector is

{\footnotesize
\begin{verbatim}
0.001242|0.888965|0.421266|-0.145791|0.225573|0.022896|0.422207|0.007678|
0.005904|-0.000997|-0.006601|-0.007472|-0.022484|-0.025189|0.080416|-0.021409
\end{verbatim}
}

\noindent and the E14 artifact vector is

{\footnotesize
\begin{verbatim}
0.007023|0.909779|0.115403|-0.249275|0.011511|0.026911|0.114251|-0.079727|
0.035886|-0.003278|-0.011850|0.062496|0.005226|-0.034832|0.027557|-0.019608
\end{verbatim}
}

\noindent again concatenating the two lines with no intervening whitespace.

\subsection{Successor choice and the forced-endgame ladder}
\label{app:params-endgame}

Two decisions in the endgame described in \S\ref{app:dialect-endgame} are made by
the agent rather than by the rules, and both belong to the \vfast{} policy rather
than to the driver.

The first is the successor choice. When a declaration leaves a player cardless
and the turn reaches them, the policy hands it to the live teammate with the
strongest publicly estimable ask: for each live teammate $u$ it scores
\[
  \max_{H} \Bigl( \Pr[u \text{ holds a card of } H] \cdot
  \max_{c \in H} \Pr[c \text{ is with an opponent}] \Bigr)
\]
over the live half-suits, using Fast per-card marginals, and gives the turn to
the maximiser. The score is computed from the public record and the policy's own
hand only, so no teammate's private holding enters it; an invalid choice is
replaced by the driver with the lowest-indexed live teammate.

The second is the willingness bit. When a whole team is cardless, the driver
sweeps the descending ladder of \S\ref{app:dialect-endgame} and the policy
answers, for each remaining half-suit at each level $\theta$, whether it would
declare that half-suit at confidence at least $\theta$. The confidence used is
$\hat\alpha$ from the same estimator that feeds \eqref{eq:stop}, and only the bit
crosses between seats. Under the \vthree{} dialect the ladder has two levels
only. Because the ladder restarts after each declaration, the earlier and safer
declarations announce exact allocations and thereby sharpen the capacity
constraints available to the ones that follow. How much of a confidence ordering
among seats this recovers is not measured, and the ordering that would reveal it
is excluded as an information leak.

\subsection{The policy-specification grammar}

Every agent in the engine is named by a policy-specification string, parsed by
\texttt{makeAgent} in \filepath{engine/src/factory.hpp}. The grammar is

\begin{center}
\texttt{name} \quad or \quad \texttt{name:key=value,key=value,\dots}
\end{center}

\noindent where \texttt{name} selects the policy family and each
\texttt{key=value} pair overrides one field of that family's configuration.
Unrecognised keys are ignored; keys are order-independent; there is no
whitespace. The example \texttt{v04:belief=block,decl=0.95,topk=6} constructs the
\vfast{} policy family with the exact reference belief substituted
(\texttt{belief=block}, i.e.\ \vblock{}), the declaration threshold raised to
$0.95$, and the two-ply lookahead width set to $K = 6$; every other field keeps
its compiled value.

The keys used in this paper are as follows. \texttt{belief} selects the inference
path and accepts \texttt{fast} (the default, \vfast{}), \texttt{block} (the exact
reference engine, \vblock{}), \texttt{exact} (the card-level capacity dynamic
program, without C5), \texttt{exactdisj}, \texttt{sinkhorn} (plain Sinkhorn,
without C5), \texttt{indep} (no capacity conditioning) and \texttt{hybrid}. The
ablation rows of \S\ref{sec:ablations} that use \texttt{exact},
\texttt{sinkhorn} and \texttt{indep} are therefore not exact-inference rows, and
are labelled accordingly there. \texttt{decl}, \texttt{lockthr},
\texttt{minteam}, \texttt{askfloor}, \texttt{pool}, \texttt{oppfloor},
\texttt{vweight}, \texttt{lweight}, \texttt{vmargin}, \texttt{ptheta},
\texttt{pphi}, \texttt{topk}, \texttt{chain} and \texttt{threat} set the
corresponding coordinates of Table~\ref{tab:bounds}. \texttt{gate} and
\texttt{mgate} set the two cheap declaration pre-gates of \S\ref{sec:declare};
\texttt{gateaudit=1} enables the instrumentation used by E16. \texttt{value},
\texttt{vdecl}, \texttt{patient}, \texttt{gmap} and \texttt{declare} are boolean
switches controlling, respectively, the one-ply value lookahead in
\eqref{eq:askutil}, the value-based stopping rule of \eqref{eq:stop}, the
patience switch, greedy maximum-a-posteriori allocation naming, and declaration
itself. \texttt{souter}, \texttt{sinner} and \texttt{particles} control the
inference solver. Individual coefficients may be set as \texttt{w0}--\texttt{w19}
and \texttt{v0}--\texttt{v15}, or in bulk as \texttt{weights=},
\texttt{vweights=} and \texttt{allparams=}.

The shipping specification used in every reported experiment is
\texttt{v04:mgate=0.008}, which is equal to the compiled defaults; it is written
explicitly in the run scripts so that the gate value appears in the recorded
command line.
